\section{Conclusion and Future Work}
\label{sec:conclusion}

This report has asked whether the register-token intervention of Darcet~\etal~\cite{darcet2024registers}, originally proposed for classification Vision Transformers, transfers to a CTC-based Handwritten Text Recognition system and whether it produces attention maps clean enough to support downstream tasks such as writer identification. The answer is qualified. The recognition half of the claim transfers cleanly: on both IAM and READ2016 the CER is flat across register counts within seed noise. The attention half transfers partially: register-added ViTs show measurably lower attention entropy, better spatial ordering, and more heads aligned with left-to-right reading, but a residual sink pattern intrinsic to CTC self-attention survives on the patch tokens and would require either an additional query-side mechanism or substantially more training data to fully suppress. The per-character attention visualization of Gurav~\etal~\cite{gurav2025beyondmem} was reproduced end-to-end from the trained model. The best absolute recognizer on IAM is the compact CNN--RNN baseline at $4.62\%$ CER, and the best interpretable recognizer is ViT-RGTS with $R = 4$ registers at $5.55\%$ CER: a $0.93\%$ CER cost for the attention behaviour the register mechanism was chosen to deliver.

One architectural finding is worth stating outside the register discussion. Two independent configurations that use a row-major 2D-patch layout collapse to roughly $73\%$ CER, isolating the strictly-left-to-right token-formation step as the single load-bearing design decision for a CTC Vision Transformer HTR encoder. This constraint applies to any variant of the family, including plain-ViT encoders that use adaptive max-pooling in place of a convolutional stem.

Four extensions follow directly from these findings.

The first is a writer-identification pipeline built on top of the ViT-RGTS attention maps. The Raven~\etal~\cite{raven2025vlac} VLAC construction aggregates character-indexed local descriptors into a per-page writer descriptor. Feeding attention-weighted patch embeddings at each character's CTC peak column into a VLAC aggregator is a direct realization of the downstream motivation of the problem statement, and it is the single most valuable follow-up.

The second is a CLS-augmented CTC ViT. Adding a single CLS token in front of the register block and supervising it with a light auxiliary loss should let the registers take over the global-information role more completely and close the residual sink pattern documented in Section~\ref{sec:res_claim}. The predicted signature is a flip of the token-norm asymmetry, with registers now carrying more norm than patches.

The third is a low-weight attention-entropy regularizer on the final-layer patch attention. CTC is agnostic to attention shape, and an explicit entropy term would provide the missing gradient signal without altering the primary decoder.

The fourth is a synthetic-to-IAM curriculum. A full pretrain on the synthetic corpus followed by real-data fine-tuning would test whether the CER gap between the from-scratch ViT-RGTS and the CNN--RNN baseline can be closed by pretraining, without altering the register-attention story.

The register-token mechanism does what its authors claim, subject to the structural differences between the classification-ViT setting it was proposed in and the CTC HTR setting it is imported into. On the accuracy side the transfer is clean. On the attention side the transfer is partial, and the mechanism of the partiality (patch tokens simultaneously acting as sink-carrying keys and CTC-decoding queries) is itself a research finding that points at a specific class of extensions. The pipeline and experiment artifacts that produced these findings are released with the report so subsequent work can build directly on top of them.
